% Options for packages loaded elsewhere
% Options for packages loaded elsewhere
\PassOptionsToPackage{unicode}{hyperref}
\PassOptionsToPackage{hyphens}{url}
%
\documentclass[
  ignorenonframetext,
]{beamer}
\newif\ifbibliography
\usepackage{pgfpages}
\setbeamertemplate{caption}[numbered]
\setbeamertemplate{caption label separator}{: }
\setbeamercolor{caption name}{fg=normal text.fg}
\beamertemplatenavigationsymbolsempty
% remove section numbering
\setbeamertemplate{part page}{
  \centering
  \begin{beamercolorbox}[sep=16pt,center]{part title}
    \usebeamerfont{part title}\insertpart\par
  \end{beamercolorbox}
}
\setbeamertemplate{section page}{
  \centering
  \begin{beamercolorbox}[sep=12pt,center]{section title}
    \usebeamerfont{section title}\insertsection\par
  \end{beamercolorbox}
}
\setbeamertemplate{subsection page}{
  \centering
  \begin{beamercolorbox}[sep=8pt,center]{subsection title}
    \usebeamerfont{subsection title}\insertsubsection\par
  \end{beamercolorbox}
}
% Prevent slide breaks in the middle of a paragraph
\widowpenalties 1 10000
\raggedbottom
\AtBeginPart{
  \frame{\partpage}
}
\AtBeginSection{
  \ifbibliography
  \else
    \frame{\sectionpage}
  \fi
}
\AtBeginSubsection{
  \frame{\subsectionpage}
}
\usepackage{iftex}
\ifPDFTeX
  \usepackage[T1]{fontenc}
  \usepackage[utf8]{inputenc}
  \usepackage{textcomp} % provide euro and other symbols
\else % if luatex or xetex
  \usepackage{unicode-math} % this also loads fontspec
  \defaultfontfeatures{Scale=MatchLowercase}
  \defaultfontfeatures[\rmfamily]{Ligatures=TeX,Scale=1}
\fi
\usepackage{lmodern}

\usetheme[]{metropolis}
\ifPDFTeX\else
  % xetex/luatex font selection
\fi
% Use upquote if available, for straight quotes in verbatim environments
\IfFileExists{upquote.sty}{\usepackage{upquote}}{}
\IfFileExists{microtype.sty}{% use microtype if available
  \usepackage[]{microtype}
  \UseMicrotypeSet[protrusion]{basicmath} % disable protrusion for tt fonts
}{}
\makeatletter
\@ifundefined{KOMAClassName}{% if non-KOMA class
  \IfFileExists{parskip.sty}{%
    \usepackage{parskip}
  }{% else
    \setlength{\parindent}{0pt}
    \setlength{\parskip}{6pt plus 2pt minus 1pt}}
}{% if KOMA class
  \KOMAoptions{parskip=half}}
\makeatother


\usepackage{longtable,booktabs,array}
\usepackage{calc} % for calculating minipage widths
\usepackage{caption}
% Make caption package work with longtable
\makeatletter
\def\fnum@table{\tablename~\thetable}
\makeatother
\usepackage{graphicx}
\makeatletter
\newsavebox\pandoc@box
\newcommand*\pandocbounded[1]{% scales image to fit in text height/width
  \sbox\pandoc@box{#1}%
  \Gscale@div\@tempa{\textheight}{\dimexpr\ht\pandoc@box+\dp\pandoc@box\relax}%
  \Gscale@div\@tempb{\linewidth}{\wd\pandoc@box}%
  \ifdim\@tempb\p@<\@tempa\p@\let\@tempa\@tempb\fi% select the smaller of both
  \ifdim\@tempa\p@<\p@\scalebox{\@tempa}{\usebox\pandoc@box}%
  \else\usebox{\pandoc@box}%
  \fi%
}
% Set default figure placement to htbp
\def\fps@figure{htbp}
\makeatother





\setlength{\emergencystretch}{3em} % prevent overfull lines

\providecommand{\tightlist}{%
  \setlength{\itemsep}{0pt}\setlength{\parskip}{0pt}}



 


\setbeamerfont{title}{size=\large} \usepackage{hyperref} \setbeamertemplate{footline}[frame number]
\makeatletter
\@ifpackageloaded{caption}{}{\usepackage{caption}}
\AtBeginDocument{%
\ifdefined\contentsname
  \renewcommand*\contentsname{Table of contents}
\else
  \newcommand\contentsname{Table of contents}
\fi
\ifdefined\listfigurename
  \renewcommand*\listfigurename{List of Figures}
\else
  \newcommand\listfigurename{List of Figures}
\fi
\ifdefined\listtablename
  \renewcommand*\listtablename{List of Tables}
\else
  \newcommand\listtablename{List of Tables}
\fi
\ifdefined\figurename
  \renewcommand*\figurename{Figure}
\else
  \newcommand\figurename{Figure}
\fi
\ifdefined\tablename
  \renewcommand*\tablename{Table}
\else
  \newcommand\tablename{Table}
\fi
}
\@ifpackageloaded{float}{}{\usepackage{float}}
\floatstyle{ruled}
\@ifundefined{c@chapter}{\newfloat{codelisting}{h}{lop}}{\newfloat{codelisting}{h}{lop}[chapter]}
\floatname{codelisting}{Listing}
\newcommand*\listoflistings{\listof{codelisting}{List of Listings}}
\makeatother
\makeatletter
\makeatother
\makeatletter
\@ifpackageloaded{caption}{}{\usepackage{caption}}
\@ifpackageloaded{subcaption}{}{\usepackage{subcaption}}
\makeatother

\usepackage{bookmark}
\IfFileExists{xurl.sty}{\usepackage{xurl}}{} % add URL line breaks if available
\urlstyle{same}
\hypersetup{
  pdfauthor={Giorgio Arcara},
  hidelinks,
  pdfcreator={LaTeX via pandoc}}


\author{Giorgio Arcara}
\date{}

\begin{document}


\begin{frame}
% TITLE SLIDES
\title{Metodologia della Ricerca Clinica\\ \vspace{1em} \emph{08 - Altri utilizzi dei test}}
\author{Giorgio Arcara,\\ Università di Padova \\ IRCCS San Camillo, Venezia}

\titlegraphic{

\vspace*{7cm}
\includegraphics[scale=0.15]{Figures/LogoSanCamilloIRCCS_Unipd_alpha.png}
\hfill \includegraphics[scale=0.15]{Figures/CC_license_3_0.png}
}
\maketitle
\end{frame}

\section{8. Altri utilizzi dei punteggi dei
test}\label{altri-utilizzi-dei-punteggi-dei-test}

\begin{frame}{Esempio}
\phantomsection\label{esempio}
\small

Testiamo un paziente tramite il test di Linguaggio Figurato 1.

Il paziente ottiene un punteggio di 12. Dopo 6 mesi viene valuato
nuovamente e ottiene un punteggio di 10. Cosa possiamo dire del
paziente?

\begin{center}
\includegraphics[width=0.5\linewidth,height=\textheight,keepaspectratio]{Figures/Figure_true_obs_1.5.png}
\end{center}

\footnotesize

Ricordiamo (vedi slides affidabilità) che in caso di ripetizioni di una
misurazione ci aspettiamo comunque delle oscillazioni nel punteggio
osservato. Il punto è capire: la variazione che ho osservato è
riconducibile a queste oscillazioni (dovute a imprecisione?) oppure ad
un cambiamento ``vero''?
\end{frame}

\begin{frame}{Valutare i cambiamenti nel tempo}
\phantomsection\label{valutare-i-cambiamenti-nel-tempo}
La valutazione dei cambiamenti nel tempo è una questione spinosa in
neuropsicologia clinica ed è molto comune.

Perché valutare lo stesso soggetto/paziente più volte?

\begin{itemize}
\item
  per indagare se è cambiato in seguito ad un trattamento di
  neuroriabilitazione.
\item
  per indagare se presenta un declino nel funzionamento cognitivo,
  possibilmente segnale di un deterioramento progressivo.
\item
  Per indagare ha subito alterazioni cognitive in seguito ad un evento,
  come un'operazione neurochirurgica per rimuovere un tumore.
\end{itemize}
\end{frame}

\begin{frame}{Valutare i cambiamenti nel tempo}
\phantomsection\label{valutare-i-cambiamenti-nel-tempo-1}
\begin{center}
  \textbf{Il problema della valutazione nel tempo}
\end{center}
\vspace{2em}

È assolutamente normale che un soggetto non ottenga esattamente lo
stesso punteggio in due diverse valutazioni.

\begin{itemize}
\item
  \emph{Effetti casuali}: legati alla precisione del test e catturati
  dall'affidabilità test-retest
\item
  \emph{Effetti sistematici}: il più comune è l'\textbf{effetto pratica}
\end{itemize}
\end{frame}

\begin{frame}{Valutare i cambiamenti nel tempo}
\phantomsection\label{valutare-i-cambiamenti-nel-tempo-2}
\begin{center}
  \textbf{L'effetto pratica}
\end{center}
\vspace{1em}

È l'effetto per cui una seconda valutazione ha un punteggio migliore
della precedente. Ci sono varie ragioni alla base dell'effetto pratica.

\begin{itemize}
\item
  Maggiore familirità con la situazione testistica (meno ansia)
\item
  Possibilità di utilizzare strategie scoperte nella prestazione
  precedente.
\end{itemize}
\end{frame}

\begin{frame}{Forme parallele}
\phantomsection\label{forme-parallele}
A volte alcuni test potrebbero avere delle \textbf{forme parallele} e
cioè una versione uguale del test in cui però gli item sono diversi.

Nelle forme parallele si introduce inoltre un'ulteriore problema è cioè
quello di dimostrare che i due test siano da considerare equivalenti.

Inoltre l'utilizzo di forme parallele non garantisce una protezione
dall'effetto pratica (Mondini e Arcara, 2008, si veda anche il test
GEMS, Mondini et al., 2022 per un esempio in cui forme parallele hanno
comunque effetto pratica).
\end{frame}

\begin{frame}{Valutare i cambiamenti nel tempo}
\phantomsection\label{valutare-i-cambiamenti-nel-tempo-3}
\begin{center}
  \textbf{Come valutare i cambiamenti nel tempo}
\end{center}
\vspace{1em}

L'unico modo (rigoroso) è utilizzare \textbf{metodi statistici per
valutare un cambiamento} (es. Chelune et al., 1993; Crawford \&
Garwhaite, 2006)

A mia conoscenza gli unici lavori in Italia che riportano tali dati non
sono molti. Ad esempio:

\begin{itemize}
\tightlist
\item
  l'\textbf{ENB-2} (Mondini et al., 2011)
\item
  \textbf{APACS} (Arcara \& Bambini, 2016)
\item
  \textbf{NADL-F} (Arcara et al., 2019)
\item
  \textbf{GEMS} (Mondini et al., 2023)
\item
  \textbf{Tele-GEMS} (Montemurro et al.~2023).
\end{itemize}
\end{frame}

\begin{frame}{Valutare i cambiamenti nel tempo}
\phantomsection\label{valutare-i-cambiamenti-nel-tempo-4}
\begin{center}
  \textbf{Come valutare i cambiamenti nel tempo}
\end{center}
\vspace{2em}

\begin{center}
\includegraphics[width=1\linewidth,height=\textheight,keepaspectratio]{Figures/thresholds_significant_changes_figlang1apacs.png}
\end{center}
\end{frame}

\begin{frame}{Valutare i cambiamenti nel tempo}
\phantomsection\label{valutare-i-cambiamenti-nel-tempo-5}
\begin{center}
  \textbf{Come valutare i cambiamenti nel tempo}
\end{center}
\vspace{2em}

\begin{center}
\includegraphics[width=1\linewidth,height=\textheight,keepaspectratio]{Figures/schermata_ENB2.png}
\end{center}
\end{frame}

\begin{frame}{Valutare i cambiamenti nel tempo}
\phantomsection\label{valutare-i-cambiamenti-nel-tempo-6}
\begin{center}
  \textbf{Metodi per valutare i cambiamenti nel tempo}
\end{center}
\vspace{2em}

Esistono numerosi metodi per valutare i cambiamenti significativi nel
tempo (si veda Collie 2002). In tutti questi metodi l'obiettivo è
costruire delle soglie (una superiore e una inferiore) suparate le quali
inferiamo che la variazione non è riconducibile al caso e, pertanto, è
da considerare come signficativa.
\end{frame}

\begin{frame}{Valutare i cambiamenti nel tempo}
\phantomsection\label{valutare-i-cambiamenti-nel-tempo-7}
\begin{center}
  \textbf{Metodi per valutare i cambiamenti nel tempo}
\end{center}
\vspace{2em}

Per farlo quasi tutti i metodi si basano su questa metodologia:

si prende un \textbf{gruppo di individui e li si testa due volte (t0 e
t1) in un lasso di tempo entro cui si assume che non sia avvenuto un
cambiamento}. Ad esempio soggetti sani testati a t0 e a t1 dopo un mese.
Questi soggetti non devono avere fatto training o altro che potrebbe
spiegare differenze tra i punteggi.

Ogni differenza tra t0 e t1 sarà dunque riconducibile a imprecisioni del
test. A partire da questo (tramite metodi statistici diversi) si
costruisce un'intervallo con un limite superiore o inferiore di
cambiamento significativo.
\end{frame}

\begin{frame}{Valutare i cambiamenti nel tempo}
\phantomsection\label{valutare-i-cambiamenti-nel-tempo-8}
\begin{center}
  \textbf{Metodi per valutare i cambiamenti nel tempo}
\end{center}
\vspace{1em}

Il principio di questi metodi è spesso analogo a quanto già visto per
identificazione di deficit/danno.

Si costruisce una sorta di distribuzione che ci dice i valori attesi
assumendo che non c'è cambiamento (H0) e si va a vedere se il punteggio
osservato (o la variazione di punteggio osservata) è così estrema da non
essere riconducibile al caso. Se è così allora si inferisce cambiamento
(peggioramento/miglioramento) significativo.
\end{frame}

\begin{frame}{Valutare i cambiamenti nel tempo}
\phantomsection\label{valutare-i-cambiamenti-nel-tempo-9}
\begin{center}
\textbf{Reliable Change Index (RCI)}
\end{center}
\vspace{2em}

L'RCI è il metodo più famoso per calcolo di differenze significative nel
tempo.

\[
  RCI = \frac{X_{i1} -X_{i0}}{SE_{diff}}
\]

\small

\[ SE_{DIFF} = \sqrt{2SE^2} \] \[SE = s_0\sqrt{1-r} \]
\[s_0 = sd(X_0)r_0 = cor(X_0,X_1)\]
\end{frame}

\begin{frame}{Valutare i cambiamenti nel tempo}
\phantomsection\label{valutare-i-cambiamenti-nel-tempo-10}
\begin{center}
\textbf{Reliable Change Index (RCI)}
\end{center}
\vspace{2em}

Una volta ottenuto l'RCI può essere trattato come uno z-score.

Se \textgreater{} 1.64, indica un miglioramento significativo (5\%),
mentre se \textless{} -1.,64, indica un peggioramento significativo
(5\%).
\end{frame}

\begin{frame}{Valutare i cambiamenti nel tempo}
\phantomsection\label{valutare-i-cambiamenti-nel-tempo-11}
\begin{center}
\textbf{Modified Reliable Change Index (mRCI)}
\end{center}
\vspace{2em}

Il mRCI è una variazione del RCI, che tiene conto di effetto pratica.

\[
mRCI = \frac{(X_{i1}-X_{i0})-m}{SE_{diff}}
\]

\small

\[m = \bar{X}_1 - \bar{X}_0\] \[SE_{diff} = \sqrt{2SE^2}\]
\[SE = s_0\sqrt{1-r_{01}}\] \[s_0 = sd(X_0)r_{01} = cor(X_0, X_1)\]
\end{frame}

\begin{frame}{Valutare i cambiamenti nel tempo}
\phantomsection\label{valutare-i-cambiamenti-nel-tempo-12}
\begin{center}
\textbf{Modified Reliable Change Index (mRCI)}
\end{center}
\vspace{2em}

Di nuovo, una volta ottenuto il mRCI può essere trattato come uno
z-score.

Se \textgreater{} 1.64, indica un miglioramento significativo (5\%),
mentre se \textless{} -1.,64, indica un peggioramento significativo
(5\%).
\end{frame}

\begin{frame}{Valutare i cambiamenti nel tempo}
\phantomsection\label{valutare-i-cambiamenti-nel-tempo-13}
\begin{center}
\textbf{Regression Method Crawford \& Garthwaithe, 2006}
\end{center}
\vspace{2em}

Anche se mRCI tiene conto di effetto pratica, non tiene conto di
regressione verso la media (i.e., punteggi estremi a t0 tenderanno a
essere più vicini alla media a t1).

Essi inoltre non tengono conto di differenze che ci potrebbero essere a
seconda del punteggio alla baseline (a t0), l'effetto pratica è corretto
sempre nella stessa maniera.
\end{frame}

\begin{frame}{Valutare i cambiamenti nel tempo}
\phantomsection\label{valutare-i-cambiamenti-nel-tempo-14}
\begin{center}
\textbf{Regression Method Crawford \& Garthwaithe, 2006}
\end{center}
\vspace{2em}

Il metodo di Crawford \& Garthwaite, tiene invece conto di tutti questi
aspetti e inoltre fornisce una statistica che tiene conto della
differenza tra campione e popolazione (e funziona anche con campioni
piccoli n\textasciitilde5).

RCI e MRCI invece non tengono conto di distinzione campione e
popolazione (come z-score).

Approfondire la formula va al di là degli scopi di questo corso.
\end{frame}

\begin{frame}{Valutare i cambiamenti nel tempo}
\phantomsection\label{valutare-i-cambiamenti-nel-tempo-15}
\begin{center}
\textbf{Come valutare i cambiamenti nel tempo}
\end{center}
\vspace{2em}

Cosa fare se il test non ha valori di cambiemanto significativo?

Fare MOLTA attenzione nell'utilizzo per inferire cambimento (evitarlo se
possibile).

Vedere i valori di affidabilità test-restest e selezionare test in cui
sono alti, ma ricordando che non è incluso l'effetto pratica.
\end{frame}

\begin{frame}{Valutare i cambiamenti nel tempo}
\phantomsection\label{valutare-i-cambiamenti-nel-tempo-16}
\begin{center}
\textbf{Da ricordare}
\end{center}
\vspace{1em}

È da ricordare che come per l'affidabilità test-retest anche gli
intervalli di cambiamento significativo si riferiscono allo specifico
intervallo temporale con cui sono stati costruiti (es. 1 mese) ed
esistono infinite distanze temporali (tra t0 e t1) per cui calcolare
intervalli di cambiamento.

In genere queste sono fatte a tempi relativamente corti (es. 1 mese)
perché altrimenti diventa difficile assumere che non c'è stato un
cambiamento.

In alternativa si usano gli stessi intervalli di rilevanza clinica. (es.
2 settimane per una operazione chirurgica, 6 mesi per Sclerosi
Multipla).
\end{frame}

\begin{frame}{Confrontare punteggi di test diversi}
\phantomsection\label{confrontare-punteggi-di-test-diversi}
\begin{center}
  \textbf{Confrontare punteggi di test diversi}
\end{center}
\vspace{2em}

In certi casi può essere utile confrontare test diversi.

Possono essere confrontati:

\begin{itemize}
\item
  test diversi che misurano lo stesso costrutto (per approfondire se la
  prestazione deficitaria è legata ad un test specifico).
\item
  test che misurano cose diverse (per inferire deficit selettivi)
\end{itemize}
\end{frame}

\begin{frame}{Confrontare punteggi di test diversi}
\phantomsection\label{confrontare-punteggi-di-test-diversi-1}
\begin{center}
  \textbf{Confrontare punteggi di test diversi}
\end{center}

Per permettere il confronto tra test diversi bisogna che i punteggi
siano trasformati in modo da essere confrontabili.

La strategia più diffusa è trasformare il punteggio sulla base della
percentuale di campione o popolazione che ha ottenuto quel punteggio (o
un punteggio inferiore; ad esempio, il 20\%, il 30\%, il 40\%, etc.).

Questo è spesso fatto tramite:

\begin{itemize}
\tightlist
\item
  Conversione in \textbf{punti--z} (ma le distribuzioni dei test devono
  essere simili)
\item
  Utilizzo dei \textbf{punteggi equivalenti} (Capitani e Laiacona, 2007)
\item
  (esistono altri metodi, es. Crawford, Howell \& Garthwaite, 1998)
\end{itemize}
\end{frame}

\begin{frame}{Confrontare punteggi di test diversi}
\phantomsection\label{confrontare-punteggi-di-test-diversi-2}
\begin{center}
  \textbf{Limiti del confronto di test diversi}
\end{center}
\vspace{2em}

È importante ricordarsi alcuni aspetti nell'utilizzo di strategie di
confronto di test diversi tramite trasformazioni.

In particolare dobbiamo ricordarci che il confronto è spesso da
considerarsi come descrittivo (e non legato ad una inferenza
statistica).
\end{frame}

\begin{frame}{Confrontare punteggi di test diversi}
\phantomsection\label{confrontare-punteggi-di-test-diversi-3}
\begin{center}
  \textbf{Limiti del confronto di test diversi}
\end{center}
\vspace{2em}

Supponiamo di traformare dei punteggi e di vedere che la persona si
colloca al 70\% nei test di memoria, mentre al 20\% per i test di
funzioni esecutive. Si potrebbe interpretare che c'è un declino
selettivo delle funzioni esecutive.

Questo è possibile (e ragionevole), ma è importante ricordare che non
sappiamo quanto \emph{improbabile} è osservare questo pattern nei sani.
Molte delle inferenze che abbiamo fatto erano basate su improbabilità,
ma in questo caso noi non sappiamo se è un pattern effettivamente molto
improbabile da suggerire un danno selettivo (lo stiamo \emph{assumendo
implicitamente}).
\end{frame}




\end{document}
